\documentclass[10pt]{article}
\usepackage[paperwidth=42cm, paperheight=30cm, margin=1.5cm]{geometry} % A3 Landscape
\usepackage[utf8]{inputenc}
\usepackage[T1]{fontenc}
\usepackage[brazil]{babel}
\usepackage{lmodern}
\renewcommand{\familydefault}{\sfdefault}

\usepackage{tikz}
\usetikzlibrary{shapes.geometric, shapes.misc, arrows.meta, positioning, calc, shadows, fit, backgrounds, spy}
\usepackage{xcolor}
\usepackage{amsmath}
\pagestyle{empty}

% Paleta de Cores Baseada no Infográfico
\definecolor{icured}{RGB}{205,32,44}       % LED Red
\definecolor{icublue}{RGB}{0,111,180}      % OxyHb / SpO2
\definecolor{icugrey}{RGB}{210,210,210}    % Fundo caixas cinzas
\definecolor{icuorange}{RGB}{240,110,0}    % IR LED / IR Signal
\definecolor{icudark}{RGB}{60,60,60}       % Texto escuro / Bordas
\definecolor{iculight}{RGB}{240,240,245}   % Tinta pastel para fundo
\definecolor{hcblack}{RGB}{30,30,30}       % Tela de Monitor

\tikzset{
    basebox/.style={
        rectangle, rounded corners=6pt, draw=icudark, thick,
        inner sep=8pt, font=\sffamily\fontsize{10}{12}\selectfont, align=left,
        text=black, drop shadow={shadow xshift=1.5pt, shadow yshift=-1.5pt}
    },
    titlebox/.style={
        basebox, fill=icugrey, text=icudark, font=\sffamily\bfseries\itshape,
        align=center
    },
    textbox/.style={
        basebox, fill=iculight, draw=icugrey, text width=6.5cm,
        font=\sffamily\fontsize{10.5}{13}\selectfont
    },
    monitor/.style={
        rectangle, rounded corners=4pt, draw=icugrey, ultra thick,
        fill=hcblack, minimum width=14cm, minimum height=6cm
    },
    wavebox/.style={
        rectangle, draw=none, fill=none, text=black, align=center,
        font=\sffamily\fontsize{9}{11}\selectfont
    },
    arrowline/.style={
        draw=icugrey!80!black, ultra thick, -{Latex[length=3mm, width=2.5mm]}
    }
}

\newcommand{\titem}[1]{\hangindent=2ex\mbox{$\bullet$ }#1\par}
\newcommand{\toptitle}[1]{\textbf{\uppercase{#1}}\\[0.5ex]}

\begin{document}
\begin{center}
\vspace*{\fill}
\resizebox{!}{0.95\textheight}{%
\begin{tikzpicture}[auto]

% ==============================
% 1. CABEÇALHO E TÍTULO
% ==============================
\node[font=\sffamily\fontsize{36}{40}\bfseries, text=black, anchor=west] (title) at (-19, 13.5) {OXIMETRIA DE PULSO};
\node[font=\sffamily\fontsize{16}{20}\selectfont, text=icudark, anchor=west] (author) at (-7.5, 13.5) {por Nick Mark MD (\textit{Tradução})};

\node[font=\sffamily\fontsize{11}{13}\selectfont, text=black, anchor=north west, text width=26cm] (principle) at (-19, 12.3) {
    \textbf{PRINCÍPIO:}\\
    A oximetria de pulso é a medição contínua e não invasiva da saturação de oxigênio da hemoglobina. Ela explora o fato de que a hemoglobina oxigenada (OxyHb) e a desoxigenada (DeoxyHb) \textbf{\textcolor{icublue}{absorvem luz vermelha e infravermelha de forma diferente}}.
};

% Caixas de Fluxo (Princípio Físico)
\node[titlebox, text width=5cm] (p1) at (-16, 10.0) {Emitir luz Vermelha (Red) e Infravermelha (IR) através da pele e medir a absorção};
\node[titlebox, text width=5cm] (p2) at (-10, 10.0) {DeoxyHb e OxyHb possuem diferentes perfis de absorção de luz};
\node[titlebox, text width=5cm] (p3) at (-4, 10.0) {O sinal que varia com o tempo representa o fluxo sanguíneo pulsátil};
\node[titlebox, text width=5cm] (p4) at (2, 10.0) {Usando uma curva padrão, a razão da absorção Red/IR (razão de modulação) calcula a SpO2};

\draw[arrowline] (p1) -- (p2);
\draw[arrowline] (p2) -- (p3);
\draw[arrowline] (p3) -- (p4);

% Explicativos acima do Monitor
\node[font=\sffamily\fontsize{10}{12}\selectfont, align=center, text width=4cm] at (8, 10.5) {A forma de onda da \textbf{pletismografia} mostra o sinal ao longo do tempo (pulso)};
\node[font=\sffamily\fontsize{10}{12}\selectfont, align=center, text width=3cm] at (12.5, 10.5) {\textbf{Amplitude} aumenta com maior Volume Sistólico};
\node[font=\sffamily\fontsize{10}{12}\selectfont, align=center, text width=3cm] at (16.5, 10.5) {\textbf{Índice de Perfusão} indica a força do sinal};

% ==============================
% 2. GRÁFICOS E ESQUEMAS (Meio)
% ==============================

% Dedo e Sensores (Esquema Simplificado)
\begin{scope}[shift={(-17, 6.5)}]
    \draw[thick, fill=orange!20, rounded corners=8pt] (-2, -0.5) rectangle (1, 0.5); % dedo base
    \draw[ultra thick] (-1, 1) -- (0, 1) -- (0.5, 0.5) -- (1, 0.5) -- (1, -0.5) -- (0, -0.5) -- (-1, -1); % Clipe
    \node[font=\sffamily\bfseries, text=icured] at (0, 1.8) {LED \textcolor{icuorange}{IR} e \textcolor{icured}{Red}};
    \node[font=\sffamily\bfseries, text=black] at (0, -1.5) {Fotodiodo};
    % Setas de luz
    \draw[ultra thick, -Latex, icured] (0, 0.8) -- (0, -0.2);
    \draw[ultra thick, -Latex, icuorange] (0.3, 0.8) -- (0.3, -0.2);
\end{scope}

% Gráfico 1: Absorção x Comprimento de Onda
\begin{scope}[shift={(-12, 4)}]
    % Eixos
    \draw[thick] (0,0) -- (5,0) node[midway, below=0.5cm, font=\sffamily\bfseries] {Comprimento de Onda (nm)};
    \draw[thick] (0,0) -- (0,4) node[midway, left=0.6cm, font=\sffamily\bfseries, rotate=90] {Absorção (\%)};
    
    % Linhas e Labels
    \draw[ultra thick, icublue] (0, 3.5) .. controls (1.5, 1.5) and (3, 2) .. (5, 1.5); % OxyHb
    \node[right, font=\sffamily\bfseries, text=icublue] at (4.2, 2.5) {OxyHb};
    
    \draw[ultra thick, icured] (0, 2.0) .. controls (1, 0.5) and (3, 2.5) .. (5, 2.5); % DeoxyHb
    \node[right, font=\sffamily\bfseries, text=icured] at (4.2, 1.8) {DeoxyHb};

    % Marcadores Red / IR
    \draw[dashed] (1, 0) -- (1, 4) node[above, font=\sffamily\bfseries, text=icured] {Red (660 nm)};
    \draw[dashed] (3.5, 0) -- (3.5, 4) node[above, font=\sffamily\bfseries, text=icuorange] {IR (940 nm)};
\end{scope}

% Gráfico 2: Razão de Modulação x SaO2
\begin{scope}[shift={(-4, 4)}]
    \draw[thick] (0,0) -- (4,0) node[midway, below=0.5cm, font=\sffamily\bfseries] {SaO2 (\%)};
    \draw[thick] (0,0) -- (0,4) node[midway, left=0.6cm, font=\sffamily\bfseries, rotate=90] {Modulation Ratio (R)};
    
    % Curva descendente
    \draw[ultra thick, icublue] (0, 3.8) .. controls (2, 3.5) and (3, 2.5) .. (4, 1.0);
    
    % Insets de onda (Red > IR vs IR > Red)
    \node[font=\sffamily\bfseries, text=icured] at (3.5, 3.8) {Red > IR};
    \node[font=\sffamily\bfseries, text=icuorange] at (3.5, 0.5) {IR > Red};
\end{scope}

% Monitor Screen
\begin{scope}[shift={(6, 6)}]
    \node[monitor] (tela) at (6, -1) {};
    
    % EKG Wave (Verde)
    \draw[thick, green!80!black] (0, -0.5) -- (0.5, -0.5) -- (0.8, -0.3) -- (1.0, -1.0) -- (1.2, 0.8) -- (1.4, -0.5) -- (1.8, -0.5) -- (2.0, -0.2) -- (2.2, -0.5) -- (4, -0.5);
    \draw[thick, green!80!black] (4, -0.5) -- (4.5, -0.5) -- (4.8, -0.3) -- (5.0, -1.0) -- (5.2, 0.8) -- (5.4, -0.5) -- (5.8, -0.5) -- (6.0, -0.2) -- (6.2, -0.5) -- (8, -0.5);
    
    % Pleth Wave (Azul claro / Ciano)
    \draw[thick, cyan!80!white] (0, -3.0) .. controls (0.5, -3.0) and (0.8, -1.5) .. (1.0, -1.5) .. controls (1.2, -1.5) and (1.2, -2.5) .. (1.5, -2.5) .. controls (1.5, -2.0) and (2.0, -3.0) .. (4, -3.0);
    \draw[thick, cyan!80!white] (4, -3.0) .. controls (4.5, -3.0) and (4.8, -1.5) .. (5.0, -1.5) .. controls (5.2, -1.5) and (5.2, -2.5) .. (5.5, -2.5) .. controls (5.5, -2.0) and (6.0, -3.0) .. (8, -3.0);
    
    % Textos da tela
    \node[font=\sffamily\bfseries\Huge, text=green!80!black] at (11, 0) {85};
    \node[font=\sffamily\bfseries, text=green!80!black] at (11, 0.8) {HR};
    
    \node[font=\sffamily\bfseries\Huge, scale=1.5, text=cyan!80!white] at (11, -2.5) {96};
    \node[font=\sffamily\bfseries, text=cyan!80!white] at (10, -2.0) {SpO2};
    \node[font=\sffamily\bfseries, text=cyan!80!white] at (12, -2.0) {PI 0.8};
\end{scope}

% ==============================
% 3. TEXTOS EM GRADE (Baixo)
% ==============================

\node[font=\sffamily\Large\bfseries, text=black, anchor=west] at (-19, 0) {PRECISÃO DA OXIMETRIA DE PULSO:};
\node[font=\sffamily\normalsize, text=black, anchor=west] at (-19, -0.7) {Vários fatores \textbf{\textcolor{icublue}{podem afetar a precisão das leituras da oximetria de pulso}}, incluindo:};

\begin{scope}[shift={(0, -2)}]

    % Linha 1
    \node[textbox, anchor=north] (t1) at (-15.5, 0) {
        \toptitle{Tipo de Sensor e Posicionamento}
        Pode ser aplicado em diferentes locais:
        \titem{\textbf{Dedo -- mais preciso}; polegar pode ser superior a outros dedos.}
        \titem{\textbf{Orelhas e Testa} -- levemente menos precisos, mas podem ser \textbf{\textcolor{icublue}{mais confiáveis em estados de vasoconstrição}} ou \underline{hipotermia}.}
        *Não há sensor ideal; tentativa e erro ajudam a encontrar a melhor posição.*
    };

    \node[textbox, anchor=north] (t2) at (-8.5, 0) {
        \toptitle{Estados de Baixo Fluxo}
        Estados de baixa perfusão (como alta vasoconstrição ou baixo débito cardíaco) podem tornar o \textbf{\textcolor{icublue}{sinal de oximetria fraco ou indetectável}}. Isso dificulta o monitoramento de SpO2 em pacientes com \textbf{fluxo não-pulsátil} (como na ECMO ou LVAD).
    };

    \node[textbox, anchor=north] (t3) at (-1.5, 0) {
        \toptitle{Menor Precisão em Baixo SpO2}
        Devido à curva de calibração desenvolvida apenas com voluntários saudáveis, a \textbf{\textcolor{icublue}{SpO2 medida pode diferir significativamente da SaO2 em valores baixos}} (ex. SpO2 menor que 75\%).
    };

    % Linha 2
    \node[textbox, anchor=north] (t4) at (-15.5, -6.5) {
        \toptitle{Cor da Pele e Esmalte}
        A oximetria pode \underline{superestimar a SpO2 em} \textbf{\textcolor{icublue}{indivíduos de pele negra}} (comparado à SaO2 arterial), principalmente em hipoxemia oculta. \textbf{Pacientes negros têm 3x mais risco de hipoxemia oculta}.
        Esmaltes de unha podem diminuir a precisão, \textbf{\textcolor{icublue}{particularmente cores escuras}} (azul, verde, preto, marrom).
    };

    \node[textbox, anchor=north] (t5) at (-8.5, -5.5) {
        \toptitle{Hemoglobina Anormal}
        \titem{\textbf{Metemoglobinemia (MetHb):} causa leitura \textbf{\textcolor{icublue}{falsa}}, cravando perto de 85-88\%.}
        \titem{\textbf{Carboxihemoglobinemia (CoHb):} causa leitura \textbf{falsamente normal} de SpO2 na casa dos 90-100\% (intoxicação por CO), escondendo a hipóxia grave.}
        \titem{\textbf{HbA1c > 7:} pode causar \textbf{\textcolor{icublue}{superestimativa da SpO2}}, embora o efeito seja pequeno.}
    };

    \node[textbox, anchor=north] (t6) at (-1.5, -4.5) {
        \toptitle{Hiperoxemia}
        A hiperóxia é prejudicial (especialmente após \underline{parada cardíaca}), mas a oximetria de pulso não distingue um PaO2 normal de um supra-normal se a SpO2 está em 100\%. O alvo é SpO2 $\ge$ 94\%.
    };

    \node[textbox, anchor=north] (t7) at (-1.5, -8.5) {
        \toptitle{Tempo de Atraso (Lag Time)}
        Demora de \textbf{5 a 15 segundos} pro sangue ir do coração até a ponta do dedo. O \textbf{atraso (lag time)} é menor se o sensor for central e maior se o débito for baixo. Por isso o SpO2 pode continuar caindo num primeiro momento até a reoxigenação chegar lá.
    };

    \node[textbox, anchor=north] (t8) at (-1.5, -13.5) {
        \toptitle{Efeitos de Medicamentos}
        Azul de metileno, indocianina verde entre outros podem causar falsas leituras \textbf{baixas de SpO2}.
    };
    
\end{scope}

% ==============================
% 4. ANÁLISE DE ONDA PLETISMOGRÁFICA (Direita Baixo)
% ==============================
\begin{scope}[shift={(10, 0)}]
    \node[font=\sffamily\Large\bfseries, text=black, anchor=west] at (-2, 0) {\uppercase{Análise da Onda Pletismográfica}};
    
    \node[font=\sffamily\normalsize, text=black, anchor=north west, text width=6cm] (descwave) at (-2, -0.5) {
        A onda pletismográfica possui componentes sistólicos \& diastólicos; o exame pode prover dicas fisiológicas sobre o \textbf{tônus e a compliance vascular}.
    };

    % Gráfico base ilustrativo
    \begin{scope}[shift={(6, -1.5)}]
        \draw[thick, gray, fill=gray!20] (0, 0) .. controls (0.5, 0) and (0.8, 1.5) .. (1.0, 1.5) .. controls (1.2, 1.5) and (1.2, 0.5) .. (1.5, 0.5) .. controls (1.5, 1.0) and (2.0, 0) .. (3, 0);
        
        \node[font=\sffamily\bfseries, text=green!70!black] at (0.2, 1.5) {Pico Sistólico};
        \node[font=\sffamily\bfseries, text=orange] at (2.2, 1.2) {Pico Diastólico};
        
        \draw[dashed] (1.0, 0) -- (1.0, 1.5);
        \draw[dashed] (1.5, 0) -- (1.5, 1.0);
        \draw[latex-latex, thick] (1.0, 1.0) -- (1.5, 1.0) node[midway, above, font=\sffamily\bfseries] {$\Delta$T};
    \end{scope}

    % Variações
    \node[wavebox, anchor=north] at (0, -4) {
        \textbf{NORMAL}\\
        \tikz{\draw[thick, icublue] (0, 0) .. controls (0.2, 0) and (0.4, 0.8) .. (0.5, 0.8) .. controls (0.6, 0.8) and (0.6, 0.3) .. (0.7, 0.3) .. controls (0.7, 0.6) and (1.0, 0) .. (1.5, 0);}
    };

    \node[wavebox, anchor=north] at (6, -4) {
        \textbf{PULSAÇÕES VENOSAS (TR Grave)}\\
        \textcolor{icublue}{Pico diastólico proeminente e tardio}\\
        \tikz{\draw[thick, icublue] (0, 0) .. controls (0.2, 0) and (0.4, 0.8) .. (0.5, 0.8) .. controls (0.6, 0.8) and (0.6, 0.2) .. (0.8, 0.2) .. controls (1.0, 1.0) and (1.2, 0) .. (1.5, 0);}
    };

    \node[wavebox, anchor=north] at (0, -7) {
        \textbf{SINAL FRACO}\\
        \textcolor{icublue}{$\downarrow \downarrow$ índice de perfusão}\\
        \tikz{\draw[thick, icublue] (0, 0) .. controls (0.2, 0) and (0.4, 0.2) .. (0.5, 0.2) .. controls (0.6, 0.2) and (0.8, 0) .. (1.5, 0);}
    };

    \node[wavebox, anchor=north] at (0, -10) {
        \textbf{VASOCONSTRIÇÃO}\\
        \textcolor{icublue}{$\uparrow$ pico diastólico}\\
        \tikz{\draw[thick, icublue] (0, 0) .. controls (0.2, 0) and (0.4, 0.8) .. (0.5, 0.8) .. controls (0.6, 0.8) and (0.6, 0.5) .. (0.7, 0.5) .. controls (0.7, 0.7) and (1.0, 0) .. (1.5, 0);}
    };

    \node[wavebox, anchor=north] at (6, -10) {
        \textbf{VASODILATAÇÃO}\\
        \textcolor{icublue}{$\uparrow$ pico sistólico, $\downarrow$ pico diastólico}\\
        \tikz{\draw[thick, icublue] (0, 0) .. controls (0.2, 0) and (0.4, 1.2) .. (0.5, 1.2) .. controls (0.6, 1.2) and (0.6, 0.1) .. (0.7, 0.1) .. controls (0.7, 0.3) and (1.0, 0) .. (1.5, 0);}
    };

    % CO-Oximetry Box
    \node[basebox, fill=icugrey, anchor=north, text width=9cm] at (3, -13) {
        \toptitle{CO-Oximetria}
        Invés de apenas 2 comprimentos de onda, os co-oxímetros utilizam \textbf{4 comprimentos de luz} para medir acuradamente hemoglobinas anormais como a CarboxiHb e MetHb.
    };

\end{scope}

\end{tikzpicture}%
}
\vspace*{\fill}
\end{center}
\end{document}
